% Integration fragment: requires amsmath, amssymb, amsthm and theorem/proposition environments.
% Labels use fcy:; bibliography keys are shared with the full article.
\section{An ordinary convergent generator for symmetric Tornheim jets}
\label{bhr:fcy:section}

The incoming \emph{Harmonic Parity and Resolvent Identities} report
already determines the quartic directional layer and the formal
all-order kernel. Its Question R10 asks for additional
analytic information about the symmetric quadratic germ controlling
cyclic fifth derivatives. We answer the stated integral version of
that question. We first give an explicit, normally convergent integral
for its full symmetric generating germ. We then extract one
transverse quadratic coefficient, including its full finite
normalization. Together with the diagonal fifth derivative this
coefficient determines every cyclic fifth derivative.

No claim that either retained coefficient has a finite evaluation in
ordinary zeta constants is made. The earlier quartic reconstruction,
formal dimension count, and cyclic factorization are used with their
existing attribution \cite{bhr:Parity}. The holomorphy of the cyclic
Tornheim sum near the origin is also a special case of Nakamura's
symmetric desingularization theorem \cite[Theorem 1.2]{bhr:Nakamura}.

\subsection{The fixed holomorphic remainder}

Write
\[
 T(A,B,C)=\sum_{m,n\ge1}\frac1{m^A n^B(m+n)^C},\qquad
 \mathfrak g(C)=\frac1{\Gamma(1+C)},\qquad \Sigma=A+B+C.
\]
The series is first used where it converges absolutely; subsequently
\(T\) denotes its meromorphic continuation. We use exactly the
normalization of the preceding report:
\begin{equation}\label{bhr:fcy:R}
 T(A,B,C)=\zeta(A)\zeta(B)
 -\frac{C\zeta(B-1)}{A+C}-\frac{C\zeta(A-1)}{B+C}
 +C R(A,B,C),
\end{equation}
where \(R\) is holomorphic and symmetric in \(A,B\) near the
origin. Here and below an apparent quotient of holomorphic germs by a
linear form is assigned its holomorphic extension when its numerator
is divisible by that form.

The finite elementary kernel needed below is
\begin{align}
 E(A,B,C)={}&\mathfrak g(C)\left\{
 \frac{\Gamma(1-A)\Gamma(1-B)}{\Sigma-2}
 +\frac{\Gamma(1-A)\zeta(B)}{A+C-1}
 +\frac{\Gamma(1-B)\zeta(A)}{B+C-1}\right\}\notag\\
 &+\zeta(A)\zeta(B)\frac{\mathfrak g(C)-1}{C}\notag\\
 &-\zeta(B-1)\frac{\Gamma(1-A)\mathfrak g(C)-1}{A+C}
 -\zeta(A-1)\frac{\Gamma(1-B)\mathfrak g(C)-1}{B+C}.
 \label{bhr:fcy:E}
\end{align}
It is jointly holomorphic near zero. For example,
\(\Gamma(1-A)\mathfrak g(-A)=1\), proving removability at \(A+C=0\).

For completeness, put \(F_A(x)=\operatorname{Li}_A(e^{-x})\) and
\begin{align}
 S(A,B;x)={}&\Gamma(1-A)\Gamma(1-B)x^{A+B-2}
 +\Gamma(1-A)\zeta(B)x^{A-1}
 +\Gamma(1-B)\zeta(A)x^{B-1}\notag\\
 &-\Gamma(1-A)\zeta(B-1)x^A
 -\Gamma(1-B)\zeta(A-1)x^B+\zeta(A)\zeta(B).
 \label{bhr:fcy:S}
\end{align}
Jonqui\`ere's expansion \cite[25.12.12]{bhr:DLMFPolylog} gives the
normally convergent integral
\begin{align}
 I_0(A,B,C)={}&\int_0^1x^{C-1}\bigl(F_A(x)F_B(x)-S(A,B;x)\bigr)\,dx
 \notag\\
 &+\int_1^\infty x^{C-1}F_A(x)F_B(x)\,dx
 \label{bhr:fcy:I0}
\end{align}
and the identity
\begin{equation}\label{bhr:fcy:RE}
 R(A,B,C)=\mathfrak g(C)I_0(A,B,C)+E(A,B,C).
\end{equation}
Indeed, integrate the six subtracted terms explicitly in the Mellin
representation for \(T\), multiply by \(\mathfrak g(C)\), and collect the
two polar terms of \eqref{bhr:fcy:R}. This gives
\eqref{bhr:fcy:E} term by term. In a sufficiently small closed parameter
polydisc the omitted Jonqui\`ere terms, after multiplication by
\(x^{C-1}\), are bounded by
\(x^{-\delta}(1+|\log x|^N)\) for some \(\delta<1\), including
any prescribed finite number of parameter derivatives. At infinity
there is uniform exponential decay. These estimates justify the
joint holomorphy and all differentiations used below.

\subsection{An explicit integral for the symmetric germ}

The earlier factorization defines the unique symmetric holomorphic
germ \(J\) by
\begin{align}
 \sum_{\mathrm{cyc}}T(A,B,C)={}&
 2\bigl(\zeta(A)\zeta(B)+\zeta(B)\zeta(C)+\zeta(C)\zeta(A)\bigr)
 \notag\\
 &-\zeta(A+B)-\zeta(B+C)-\zeta(C+A)\notag\\
 &+2\bigl(\zeta(A)+\zeta(B)+\zeta(C)\bigr)+1+ABCJ(A,B,C).
 \label{bhr:fcy:J}
\end{align}
We retain this normalization exactly. The next formula expresses
\(J\) through ordinary polylogarithms, Gamma and zeta functions,
and ordinarily convergent integrals, without any Tornheim derivative
on the right-hand side.

For a holomorphic function \(H\), define
\[
 \mathcal D_{A,B}H(A,B,C)=
 \frac{H(A,B,C)-H(A,0,C)-H(0,B,C)+H(0,0,C)}{AB}.
\]
At zero parameters this means the corresponding derivative. Set
\begin{align*}
 d_A(x)&=\frac{\operatorname{Li}_A(e^{-x})-\operatorname{Li}_0(e^{-x})}{A},&
 g_A(x)&=\frac{\Gamma(1-A)x^A-1}{A},\\
 p_A&=\frac{\zeta(A)+1/2}{A},&
 q_A&=\frac{\zeta(A-1)+1/12}{A},
\end{align*}
with the holomorphic values at \(A=0\). Define
\begin{align}
 S_\Delta(A,B;x)={}&\frac{g_A(x)g_B(x)}{x^2}
 +\frac{g_A(x)p_B+g_B(x)p_A}{x}\notag\\
 &-g_A(x)q_B-g_B(x)q_A+p_Ap_B,
 \label{bhr:fcy:SDelta}\\
 I_\Delta(A,B,C)={}&\int_0^1x^{C-1}
 \bigl(d_A(x)d_B(x)-S_\Delta(A,B;x)\bigr)\,dx\notag\\
 &+\int_1^\infty x^{C-1}d_A(x)d_B(x)\,dx.
 \label{bhr:fcy:IDelta}
\end{align}

\begin{theorem}[Ordinary convergent symmetric generator]
\label{bhr:fcy:generator}
The integrals \eqref{bhr:fcy:IDelta} converge normally in a common
neighborhood of \((0,0,0)\), including after arbitrary fixed
parameter derivatives, and
\begin{equation}\label{bhr:fcy:Jintegral}
 J(A,B,C)=\sum_{\mathrm{cyc}}
 \left\{\mathfrak g(C)I_\Delta(A,B,C)+\mathcal D_{A,B}E(A,B,C)\right\}.
\end{equation}
\end{theorem}

\begin{proof}
Apply the rectangular difference in \(A,B\) to
\eqref{bhr:fcy:I0}. The difference of the product \(F_AF_B\),
divided by \(AB\), is \(d_Ad_B\); the corresponding difference
of \(S\) is exactly \(S_\Delta\). Cauchy estimates on a slightly
larger parameter polydisc preserve the integrable bounds of
\eqref{bhr:fcy:I0}; equivalently, one can use
\(d_A(x)=g_A(x)/x+p_A-q_Ax+O(x^{2-\epsilon})\), locally
uniformly with logarithmic factors after differentiation. Thus
\(I_\Delta=\mathcal D_{A,B}I_0\), with normal convergence.

In the cyclic sum of \eqref{bhr:fcy:R}, the two terms with
denominator \(A+C\) combine to \(-\zeta(B-1)\), and similarly
for the other denominators. All remaining explicit zeta terms
depend on at most two variables. The complete rectangular
three-variable difference therefore annihilates them and gives
\[
 ABCJ(A,B,C)=\sum_{\mathrm{cyc}}C\bigl(
 R(A,B,C)-R(A,0,C)-R(0,B,C)+R(0,0,C)\bigr).
\]
To check the left-hand side, apply the same difference to
\eqref{bhr:fcy:J}: all its explicit terms vanish, and \(ABCJ\)
vanishes on every coordinate plane. Divide by \(ABC\), substitute
\eqref{bhr:fcy:RE}, and use holomorphic extension at zero parameters.
This proves \eqref{bhr:fcy:Jintegral}.
\end{proof}

The theorem is an analytic representation of the previously defined
germ. It supplies all symmetric jet coefficients in one fixed,
ordinary convergent kernel. Its content includes the endpoint
subtractions and their finite normalization; replacing the kernel by
an unsubtracted integral would not preserve its Taylor coefficients.

\subsection{The first transverse quadratic coefficient}

Write \(\gamma\) for Euler's constant and put
\[
 L=\log(2\pi),\qquad z_j=\zeta^{(j)}(0),\qquad
 y_j=\zeta^{(j)}(-1),\qquad \kappa_2=\zeta(2),\quad \kappa_4=\zeta(4).
\]
The quadratic homogeneous part of \(J\) has the form
\begin{equation}\label{bhr:fcy:J2form}
 J_2(A,B,C)=j_{20}(A^2+B^2+C^2)+\chi(AB+BC+CA).
\end{equation}
Thus \(\chi=[AB]J\) is the additional scalar to be determined.
For \(x>0\), let
\[
 f_j(x)=\left.\partial_s^j\operatorname{Li}_s(e^{-x})\right|_{s=0},
 \qquad h(x)=\log x+\gamma,
\]
and define
\begin{align}
 P_{21}(x)={}&\frac{h(h^2+\kappa_2)}{x^2}
 +\frac{z_1(h^2+\kappa_2)+z_2h}{x}
 -y_1(h^2+\kappa_2)-y_2h+z_1z_2,
 \label{bhr:fcy:P21}\\
 P_{22}(x)={}&\frac{(h^2+\kappa_2)^2}{x^2}
 +\frac{2z_2(h^2+\kappa_2)}{x}-2y_2(h^2+\kappa_2)+z_2^2,
 \label{bhr:fcy:P22}\\
 \mathcal K(x)={}&\frac14 f_2(x)^2+h(x)f_2(x)f_1(x),\qquad
 P(x)=\frac14P_{22}(x)+h(x)P_{21}(x).
 \label{bhr:fcy:KP}
\end{align}
Let
\begin{equation}\label{bhr:fcy:chiintegral}
 \mathcal I=\int_0^1\frac{\mathcal K(x)-P(x)}x\,dx
                   +\int_1^\infty\frac{\mathcal K(x)}x\,dx.
\end{equation}
Both integrals converge ordinarily and absolutely.

\begin{theorem}[Explicit transverse fifth-order coordinate]
\label{bhr:fcy:chi}
The coefficient in \eqref{bhr:fcy:J2form} is
\begin{equation}\label{bhr:fcy:chiresult}
 \chi=\mathcal I+\mathcal E,
\end{equation}
where
\begin{align}
 \mathcal E={}&-\frac{10\gamma^4+20\gamma^3+12\gamma^2\kappa_2
 +30\gamma^2+12\gamma \kappa_2+30\gamma+2\kappa_2^2+6\kappa_2+15}{16}
 \notag\\
 &+\frac{-\gamma^4-2\gamma^2\kappa_2+\kappa_2^2+2\kappa_4}{4}\,y_1
 -\frac{\gamma(\gamma^2+\kappa_2)}2\,y_2\notag\\
 &-\bigl(\gamma^3+3\gamma^2+6\gamma+6+(\gamma+1)\kappa_2\bigr)z_1
 -\frac{3\gamma^2+6\gamma+6+\kappa_2}{2}\,z_2\notag\\
 &+\frac{\gamma^2-\kappa_2}{2}\,z_1z_2+\frac\gamma4z_2^2.
 \label{bhr:fcy:chicorrection}
\end{align}
\end{theorem}

\begin{proof}
All derivatives of \(R\) in this proof are evaluated at the
origin. The coefficient of \(A^2B^2C\) in
\(\sum_{\mathrm{cyc}}CR(A,B,C)\) is
\[
 \frac14R_{AABB}+R_{AABC}.
\]
The first cyclic term contributes \(R_{AABB}/4\), and the
other two contribute \(R_{AABC}/2\) each. The explicit
zeta terms in \eqref{bhr:fcy:J} and the cyclic form of
\eqref{bhr:fcy:R} have no monomial involving all three variables.
Consequently
\begin{equation}\label{bhr:fcy:chiderivative}
 \chi=\frac14R_{AABB}+R_{AABC}.
\end{equation}

Differentiate \eqref{bhr:fcy:RE}, using \(\mathfrak g(0)=1\) and
\(\mathfrak g'(0)=\gamma\). The integral contribution to
\eqref{bhr:fcy:chiderivative} is precisely \eqref{bhr:fcy:chiintegral}:
\(S_{AAB}=P_{21}\), \(S_{AABB}=P_{22}\), and the derivative
of \(\mathfrak g(C)x^{C-1}\) at zero is \(h(x)/x\).
For a direct convergence check, the differentiated Jonqui\`ere
expansions are
\[
 f_1(x)=\frac h x+\sum_{k\ge0}
       \frac{(-1)^k\zeta'(-k)}{k!}x^k,\qquad
 f_2(x)=\frac{h^2+\kappa_2}x+\sum_{k\ge0}
       \frac{(-1)^k\zeta''(-k)}{k!}x^k.
\]
The displayed subtractions remove both singular products, the
cross terms of analytic degrees zero and one, and the constant
analytic product. Hence \((\mathcal K-P)/x=O(1+|\log x|^3)\) at zero.
At infinity \(f_1,f_2=O(e^{-2x})\), so
\(\mathcal K/x=O((1+\log x)e^{-4x}/x)\).

The remaining finite contribution is
\[
 \mathcal E=[A^2B^2]E+2[A^2BC]E.
\]
Insert
\begin{align*}
 \log\Gamma(1-A)&=\gamma A+\frac{\kappa_2}{2} A^2
                  +\frac{\zeta(3)}3A^3+\frac{\kappa_4}4A^4+O(A^5),\\
 \log \mathfrak g(C)&=\gamma C-\frac{\kappa_2}{2} C^2
                  +\frac{\zeta(3)}3C^3-\frac{\kappa_4}4C^4+O(C^5)
\end{align*}
into the elementary expression \eqref{bhr:fcy:E}. Expand its three
nonresonant denominators geometrically and divide the three
removable quotients as holomorphic power series. Extraction of
the two coefficients gives \eqref{bhr:fcy:chicorrection}. This is
finite polynomial algebra; the accompanying exact script performs
it independently with all zeta derivatives left symbolic.
\end{proof}

\subsection{Every cyclic fifth derivative}

Put \(\Omega_5=W_{1,1,1}^{(5)}(0)\), where
\(W_{a,b,c}(t)=T(at,bt,ct)\) is restricted before taking its
one-variable continuation at zero. Write
\[
 e_1=a+b+c,\quad e_2=ab+bc+ca,\quad e_3=abc,\quad
 D=e_1^2-3e_2,\quad p_2=e_1^2-2e_2.
\]

\begin{theorem}[Cyclic fifth-order formula]\label{bhr:fcy:fifth}
If the three pairwise sums of \(a,b,c\) are nonzero, then
\begin{align}
 \sum_{\mathrm{cyc}}W_{a,b,c}^{(5)}(0)={}&
 e_3p_2\Omega_5-120e_3D\chi
 -5(e_1e_2-3e_3)D\,Lz_4\notag\\
 &+20(e_1e_2^2-4e_1^2e_3+3e_2e_3)z_2z_3\notag\\
 &+(-2e_1^5+5e_1^3e_2-5e_1e_2^2
                         +47e_1^2e_3-69e_2e_3)z_5.
 \label{bhr:fcy:fifthformula}
\end{align}
For arbitrary complex slopes the same polynomial gives the fifth
derivative of the holomorphic cyclic sum restricted as a whole.
This last assertion does not assign values to its individual
singular summands.
\end{theorem}

\begin{proof}
On the diagonal, differentiation of \eqref{bhr:fcy:J} gives
\[
 j_{20}+\chi=\frac{
 \Omega_5+10Lz_4-40z_2z_3+32z_5}{120}.
\]
Thus \eqref{bhr:fcy:J2form} is completely determined by
\(\Omega_5\) and the ordinarily convergent coordinate
\eqref{bhr:fcy:chiresult}. For general slopes, the fifth derivative
of the explicit part of \eqref{bhr:fcy:J} is
\[
 -5Lz_4\sum_{\mathrm{sym}}a^4b
 +20z_2z_3\sum_{\mathrm{sym}}a^3b^2
 -\bigl((a+b)^5+(b+c)^5+(c+a)^5\bigr)z_5.
\]
The remaining derivative is \(120abcJ_2(a,b,c)\).
Substitute the preceding value of \(j_{20}\), use the elementary
symmetric polynomial identities, and collect terms. This yields
\eqref{bhr:fcy:fifthformula}. Since \eqref{bhr:fcy:J} is jointly
holomorphic, its cyclic restriction exists even when an individual
pairwise divisor would prevent restricting a single Tornheim term.
\end{proof}

For instance, the explicit asymmetric identity is
\begin{align}
 &W_{1,2,3}^{(5)}(0)+W_{2,3,1}^{(5)}(0)+W_{3,1,2}^{(5)}(0)
       -84\Omega_5\notag\\
 &\hspace{18mm}=-2160\chi-720L\zeta^{(4)}(0)
       +1200\zeta''(0)\zeta'''(0)-1704\zeta^{(5)}(0).
 \label{bhr:fcy:example}
\end{align}
For real slopes, \(D=\tfrac12((a-b)^2+(b-c)^2+(c-a)^2)\).
Consequently a single real cyclic fifth ray with \(abc\ne0\)
eliminates the displayed transverse coefficient only on the
diagonal. For complex slopes the cone \(D=0\) has additional
points. These are coefficient statements in the specified two
coordinate representation, not claims of arithmetic independence.

\needspace{13\baselineskip}
\subsection{Numerical checks and the remaining arithmetic problem}
The package evaluates \(\mathcal I\) using differentiated
Jonqui\`ere series on \((0,1)\) and exponentially convergent
incomplete-Gamma sums on \((1,\infty)\). At 70 decimal digits
of working precision the diagnostic values are
\begin{align*}
 \mathcal I&=-0.0128358838132664551878307355997471413775071472849\ldots,\\
 \mathcal E&=18.6522612153652911636439613245591942979227651526949\ldots,\\
 \chi&=18.6394253315520247084561305889594471565452580054100\ldots,\\
 \Omega_5&=8808.87952093843854353140320772504731742914386526227\ldots.
\end{align*}
An independently evaluated Mellin continuation of individual
\(T\) rays checks \eqref{bhr:fcy:fifthformula} on the cycles generated
by \((1,2,3)\), \((1,1,2)\), and \((1,-2,3)\); the largest
recorded discrepancy is below \(6\cdot10^{-30}\). These are
uncertified numerical diagnostics, not interval enclosures or
proofs of the analytic identities. The exact coefficient script
checks the elementary correction and the full symbolic fifth-order
formula.

The new generator answers the ordinary-integral alternative of
Question R10 and makes all higher symmetric coefficients concrete.
Finite ordinary-constant evaluations of \(\chi\) and
\(\Omega_5\) remain open. A sharper next problem is to find an
independent Gamma or Barnes representation of the kernel in
\eqref{bhr:fcy:chiintegral}, with its normalization \eqref{bhr:fcy:chicorrection}
unchanged. At the following order, the same generator yields the
three coefficients of the symmetric cubic germ \(J_3\); any
reduction below that formal count requires an additional identity.
